% =====================================================================
%  摘要 + 关键词（占位，最终成文时替换）
%  要求：概括全文针对 B 题四个子问题建立的模型、方法与主要结果；
%        无需翻译成英文；原则上不超过一页。
% =====================================================================
【待填充：全文摘要，建议 300--500 字。写作提纲（可按此展开）：
  (1) 背景与目标——无线电干扰源的快速自动定位与清除；
  (2) 针对问题一，建立示向度交会定位模型，……；
  (3) 针对问题二，建立第二检测点最优选择模型，……；
  (4) 针对问题三，建立多源快速定位与清除策略，……；
  (5) 针对问题四，建立定向干扰源定位模型，……；
  (6) 主要结论与数值/实测指标（清除比例、平均定位清除时间等）。】

本文针对无线电干扰源的快速自动定位与清除问题，\ldots{}（此处填写全文摘要）。

\keywords{示向度交会\quad 定位误差\quad 信息驱动搜索\quad 路径规划\quad 定向干扰源}
